\documentclass[11pt,a4paper]{article}
\usepackage[utf8]{inputenc}
\usepackage[T1]{fontenc}
\usepackage[french]{babel}
\usepackage{geometry}
\usepackage{amsmath}
\geometry{margin=2cm}

% Définition d'un environnement pour les algorithmes
\newcounter{algocounter}
\newenvironment{algorithm}[1]{
  \refstepcounter{algocounter}
  \begin{center}
  \fbox{
    \begin{minipage}{0.9\textwidth}
      \textbf{Algorithme \thealgocounter : #1}
      \vspace{0.2cm}
      \hrule
      \vspace{0.3cm}
}{
    \end{minipage}
  }
  \end{center}
  \vspace{0.3cm}
}

% Commandes pour le pseudo-code
\newcommand{\KW}[1]{\textbf{#1}}
\newcommand{\COMMENT}[1]{\textit{// #1}}
\newcommand{\INDENT}{\quad}
\newcommand{\FUNCTION}[2]{\KW{Fonction} #1(#2)}
\newcommand{\PROCEDURE}[2]{\KW{Procédure} #1(#2)}
\newcommand{\IF}[1]{\KW{Si} #1 \KW{alors}}
\newcommand{\ELSE}{\KW{sinon}}
\newcommand{\ELSEIF}[1]{\KW{sinon si} #1 \KW{alors}}
\newcommand{\ENDIF}{\KW{fin si}}
\newcommand{\FOR}[1]{\KW{Pour} #1 \KW{faire}}
\newcommand{\ENDFOR}{\KW{fin pour}}
\newcommand{\WHILE}[1]{\KW{Tant que} #1 \KW{faire}}
\newcommand{\ENDWHILE}{\KW{fin tant que}}
\newcommand{\RETURN}[1]{\KW{retourner} #1}
\newcommand{\STATE}[1]{#1}

\title{Modèles d'Algorithmes en Pseudo-Code}
\author{Notes d'Algorithmique}
\date{\today}

\begin{document}

\maketitle

\section{Introduction au Pseudo-Code}

Le pseudo-code est une méthode de description d'algorithmes qui utilise un langage simple et structuré, indépendant de tout langage de programmation spécifique.

\section{Structures de Base}

\subsection{Structure Conditionnelle}

\begin{algorithm}{Structure Si-Alors-Sinon}
\IF{condition} \\
\INDENT \STATE{instruction1} \\
\INDENT \STATE{instruction2} \\
\ELSE \\
\INDENT \STATE{instruction3} \\
\ENDIF
\end{algorithm}

\subsection{Boucle Pour}

\begin{algorithm}{Boucle Pour avec compteur}
\FOR{i de 1 à n} \\
\INDENT \STATE{instruction} \\
\ENDFOR
\end{algorithm}

\subsection{Boucle Tant Que}

\begin{algorithm}{Boucle Tant Que}
\WHILE{condition} \\
\INDENT \STATE{instruction} \\
\ENDWHILE
\end{algorithm}

\section{Exemples d'Algorithmes}

\subsection{Recherche du Maximum}

\begin{algorithm}{Recherche du maximum dans un tableau}
\FUNCTION{MaxTableau}{T, n} \\
\INDENT max $\leftarrow$ T[1] \\
\INDENT \FOR{i de 2 à n} \\
\INDENT \INDENT \IF{T[i] > max} \\
\INDENT \INDENT \INDENT max $\leftarrow$ T[i] \\
\INDENT \INDENT \ENDIF \\
\INDENT \ENDFOR \\
\INDENT \RETURN{max} \\
\KW{fin fonction}
\end{algorithm}

\subsection{Tri par Sélection}

\begin{algorithm}{Tri par sélection}
\PROCEDURE{TriSelection}{T, n} \\
\INDENT \FOR{i de 1 à n-1} \\
\INDENT \INDENT min $\leftarrow$ i \\
\INDENT \INDENT \FOR{j de i+1 à n} \\
\INDENT \INDENT \INDENT \IF{T[j] < T[min]} \\
\INDENT \INDENT \INDENT \INDENT min $\leftarrow$ j \\
\INDENT \INDENT \INDENT \ENDIF \\
\INDENT \INDENT \ENDFOR \\
\INDENT \INDENT échanger T[i] et T[min] \\
\INDENT \ENDFOR \\
\KW{fin procédure}
\end{algorithm}

\subsection{Recherche Dichotomique}

\begin{algorithm}{Recherche dichotomique dans un tableau trié}
\FUNCTION{RechercheDichotomique}{T, n, valeur} \\
\INDENT début $\leftarrow$ 1 \\
\INDENT fin $\leftarrow$ n \\
\INDENT \WHILE{début $\leq$ fin} \\
\INDENT \INDENT milieu $\leftarrow$ (début + fin) / 2 \\
\INDENT \INDENT \IF{T[milieu] = valeur} \\
\INDENT \INDENT \INDENT \RETURN{milieu} \\
\INDENT \INDENT \ELSEIF{T[milieu] < valeur} \\
\INDENT \INDENT \INDENT début $\leftarrow$ milieu + 1 \\
\INDENT \INDENT \ELSE \\
\INDENT \INDENT \INDENT fin $\leftarrow$ milieu - 1 \\
\INDENT \INDENT \ENDIF \\
\INDENT \ENDWHILE \\
\INDENT \RETURN{-1} \COMMENT{élément non trouvé} \\
\KW{fin fonction}
\end{algorithm}

\section{Conseils pour écrire du Pseudo-Code}

\begin{itemize}
    \item Utilisez des noms de variables explicites
    \item Indentez correctement pour montrer la structure
    \item Utilisez des commentaires pour clarifier les étapes complexes
    \item Soyez précis dans les conditions et les opérations
    \item Testez mentalement votre algorithme avec des exemples
\end{itemize}

\section{Template Vide}

\begin{algorithm}{Votre algorithme}
\FUNCTION{NomFonction}{paramètres} \\
\INDENT \COMMENT{Votre code ici} \\
\INDENT \\
\KW{fin fonction}
\end{algorithm}

\end{document}